\section{Evaluation}\label{sec:eval}
The following numbers are synthetic teaching data, not results from an
implemented system. Each row represents one 10-second trace with 100 arrivals,
one worker, the same fixed model, and a warm model cache. All admitted work
finishes before the trace ends. Latency is measured from arrival to completion
as in Equation~\ref{eq:latency}, over completed requests only. Rejected requests
are excluded from latency quantiles. The completion rate divides completed
requests by the 10-second trace duration. Both policies use the same estimates;
FIFO does not use them to order requests.

\begin{table}[h]
\centering
\begin{tabular}{lrrr}
Policy & Admitted/completed & P99 (ms) & Completed/s\\
FIFO & 100/100 & 100 & 10\\
GateQueue & 80/80 & 80 & 8\\
\end{tabular}
\caption{Synthetic burst trace; both configurations admit all offered work.}
\label{tab:burst}
\end{table}

At low load, both policies complete 100 of 100 requests and have P99 of 20 ms.
No latency gain appears in this low-load example. Per-request logs and repeated
traces are not included. Table~\ref{tab:burst} supports a 20\% reduction in
completed-request P99 under the respective admission policies.
